\documentclass[11pt]{article}
\usepackage[a4paper,margin=28mm]{geometry}
\usepackage{amsmath,amssymb,amsthm,booktabs,longtable,hyperref}
\newtheorem{lemma}{Lemma}
\newtheorem{theorem}[lemma]{Theorem}
\newtheorem{corollary}[lemma]{Corollary}
\theoremstyle{remark}\newtheorem{remark}[lemma]{Remark}
\newcommand{\Inv}{\operatorname{Inv}}
\newcommand{\gen}[1]{\langle #1\rangle}
\newcommand{\GL}{\operatorname{GL}}
\title{Small sets of involutions exclude parabolic overgroups}
\author{A conditional application to the existing finite intake}
\date{}
\begin{document}
\maketitle
\begin{abstract}
Let $D=\gen{a,b}$ be one of the admitted prefixes in the acquired
$G_2(11)$ model. Put $U=C_G(D)$ and $S=C_U(e)$, where $e$ is an
involution of $U$. If $U$ has fewer than five nonidentity involutions
and $5$ does not divide $|S|$, orbit--stabilizer implies that $5$ does
not divide $|U|$. When $11$ does not divide $|D|$, an embedded Levi
centre would give the opposite conclusion if $D$ lay in either maximal
parabolic. This proves the required restriction without assuming
$U=S$. We give the proof, its exact imported hypotheses, and its
application to the recorded 81 eligible prefixes. One application,
job 24, was already proved by a stronger centralizer equality; the
remaining 80 are distinct applications relative to that retained record.
No disposition of a final quartet is made here.
\end{abstract}

\section{The centralizer of one involution}

For a finite group $V$, write
\[
 \Inv(V)=\{x\in V:x^2=1,\ x\ne1\}.
\]
In particular the identity is excluded from every count below. All
centralizers are full centralizers in the indicated group.

\begin{lemma}\label{orbit}
Let $U$ be a finite group, let $e\in\Inv(U)$, and put $S=C_U(e)$ and
$N=|\Inv(U)|$. Then
\[
 [U:S]=|e^U|\leq N.
\]
If $\ell$ is a prime, $N<\ell$, and $\ell\nmid |S|$, then
$\ell\nmid |U|$.
\end{lemma}
\begin{proof}
Use the conjugation action $e^u=u^{-1}eu$. Its stabilizer is exactly
$C_U(e)=S$. Every conjugate of $e$ is again a nonidentity involution
of $U$, so orbit--stabilizer gives the displayed equality and inequality.
Set $r=[U:S]$. It is a positive integer with $r\leq N<\ell$, and
therefore $\ell\nmid r$. The equality $|U|=r|S|$ and the hypothesis
$\ell\nmid |S|$ give the conclusion.
\end{proof}

Neither the order of $U$ nor the equality $U=S$ has to be acquired.
It is important that $S$ be the \emph{full} centralizer of $e$ in $U$:
a known subgroup of that centralizer supplies neither the index nor
the factorization used in the proof.

\begin{remark}
The two numerical hypotheses have distinct roles. For example,
$U=C_2\times C_\ell$, with $\ell$ odd and $e$ the involution, has
$N=1<\ell$ but $S=U$, so the hypothesis on $|S|$ cannot be omitted.
In a dihedral group of order $2\ell$, a reflection has centralizer of
order two and there are $\ell$ involutions; this shows why the strict
inequality $N<\ell$ matters. Finally $U=A_4$, with $e$ a double
transposition, has $N=3$ and $|S|=4$, although $U\ne S$.
This last example concerns the abstract counting lemma; it is not a
claim that such a group occurs as one of the present prefix centralizers.
\end{remark}

\section{The complement argument}

We reuse the following standard part of Schur--Zassenhaus. The proof
is included to specify the conjugating subgroup and to avoid treating
an abstract Levi quotient as an embedded complement.

\begin{lemma}\label{complement}
Let $Q$ be a finite $p$-group, let $B$ have order prime to $p$, and
let $E=Q\rtimes B$. Every complement to $Q$ in $E$ is conjugate to
$B$ by an element of $Q$.
\end{lemma}
\begin{proof}
Induct on $|Q|$. There is nothing to prove if $Q=1$.
Otherwise $A=\Omega_1(Z(Q))$ is a nontrivial elementary abelian
$p$-group, characteristic in $Q$ and normal in $E$.
If $C$ is a complement to $Q$, the images of $C$ and $B$ in $E/A$
are complements to $Q/A$. Induction gives a conjugation by an element
of $Q/A$ making these images equal. Lift that element to $Q$ and
conjugate $C$. We may now suppose that $CA=BA$.
Thus $C$ and $B$ are complements to $A$ in $A\rtimes B$.

Write $A$ additively. There is a unique function $f:B\to A$ with
$C=\{f(b)b:b\in B\}$. Closure is equivalent to
\[
 f(bd)=f(b)+b\cdot f(d).
\]
Put $s=\sum_{d\in B}f(d)$ and $v=|B|^{-1}s$, where the inverse
is taken in $\mathbb F_p$. Summing the displayed equation over $d$
yields $s=|B|f(b)+b\cdot s$, whence
$f(b)=v-b\cdot v$. It follows that $C=vBv^{-1}$.
Both conjugations used in the induction lie in $Q$.
\end{proof}

\begin{corollary}\label{levi}
Suppose that $P=Q\rtimes M$ is a subgroup of a finite group $G$,
where $Q$ is a $p$-group. If $D\leq P$ has order prime to $p$,
then $C_G(D)$ contains a conjugate of $Z(M)$.
\end{corollary}
\begin{proof}
Let $\pi:P\to M$ be the projection supplied by the embedded complement
and set $B=\pi(D)$. Since $D\cap Q=1$, projection restricts to an
isomorphism $D\cong B$. In $\pi^{-1}(B)=Q\rtimes B$, both $D$
and the standard copy of $B$ in $M$ are complements to $Q$.
Lemma~\ref{complement} conjugates $D$ into $M$ by an element of $Q$.
The centre of $M$ centralizes that conjugate of $D$; conjugating back
places a conjugate of $Z(M)$ in $C_G(D)$.
\end{proof}

\section{A prime centralizer obstruction}

\begin{theorem}\label{general}
Let $a,b,e\in G$, where $G$ is finite, and suppose that $e$ is a
nonidentity involution commuting with $a$ and $b$. Put
\[
 D=\gen{a,b},\quad U=C_G(D),\quad
 S=C_G(a,b,e)=C_U(e),\quad N=|\Inv(U)|.
\]
Let $p$ and $\ell$ be primes. Assume
\[
 p\nmid |D|,\qquad N<\ell,\qquad \ell\nmid |S|.
\]
Then $D$ is contained in no subgroup $P=Q\rtimes M$ of $G$ for
which $Q$ is a $p$-group and $\ell\mid |Z(M)|$.
\end{theorem}
\begin{proof}
The equality $C_G(D)=C_G(a,b)$ follows from $D=\gen{a,b}$.
The commutations put $e$ in $\Inv(U)$, and intersection of the three
centralizer conditions gives $C_U(e)=C_G(a,b,e)=S$.
Lemma~\ref{orbit} now gives $\ell\nmid |U|$.
If $D\leq Q\rtimes M$ as in the conclusion, Corollary~\ref{levi}
puts a conjugate of $Z(M)$ inside $U$. Its order is divisible by
$\ell$, contradicting Lagrange's theorem.
\end{proof}

The theorem places no simplicity or centrelessness assumption on $G$.
Those assumptions may be needed for later quartet deductions; they are
not needed for this prefix restriction.

\section{Specialization to the admitted $G_2(11)$ model}

We use the same full named model, faithful frame, original ordered marks,
membership, original index and catalogue identifications as in the
retained prefix proof~\cite{old}. In particular, identifying the acquired
matrices with the full finite split group $G_2(11)$ is an explicit
borrowed premise. It is not deduced from the present involution counts.
For each prefix the original full stabilizer premise identifies its
recorded $S$ with $C_G(a,b,e)$, not just with a subgroup of that centralizer.

For each of the two maximal parabolic types the retained split Levi
structure is
\[
 P=Q\rtimes M,\qquad Q\text{ an }11\text{-group},\qquad
 M\cong\GL_2(11).
\]
Here $M$ is an actual embedded finite complement. The scalar centre of
$\GL_2(11)$ has order ten, so $5\mid |Z(M)|$.
The two algebraic Levi identifications and decompositions are recorded
in Lin~\cite{Lin}, Sections 7.1.4 and 7.1.6, printed pages 450--451;
the finite specialization and the distinction from a projective quotient
are retained explicitly in~\cite{old}. We reuse that imported structure.
No new verification of the primary representation or root data is
claimed here. The two conjugacy types have the same relevant structure.

\begin{corollary}\label{q11}
For any admitted prefix satisfying
\[
 N=|\Inv(C_G(a,b))|<5,\qquad
 5\nmid |C_G(a,b,e)|,\qquad 11\nmid |\gen{a,b}|,
\]
neither maximal parabolic type contains $\gen{a,b}$.
Consequently any subgroup containing that prefix, including a later
left end or a whole quartet, cannot be contained in either type.
\end{corollary}
\begin{proof}
Apply Theorem~\ref{general} with $p=11$ and $\ell=5$ to each of
the two embedded finite Levi decompositions. The proof uses only the
factor five in the Levi centre, not an exact order for $U$.
\end{proof}

\section{The exact finite application}

The adopted prefix facts~\cite{facts} contain 157 records. Their
acquisition scans every original nonidentity involution of $C_G(a)$
in its adopted bank and tests commutation with the actual $b$.
Bank completeness, group membership, involution and $a$-commutation,
the faithful model and the full original $S$ semantics remain imports.
The same acquired record supplies the exact order $m=|ab|$ by an
initial tested annihilator and final tests for every prime divisor.
Distinct noncommuting involutions then give $|D|=2m$.
This note uses the accepted records; it does not independently replay
those matrix equations.

Rechecking only the integer predicates in Corollary~\ref{q11} gives
81 records, partitioned as follows.
\begin{center}
\begin{tabular}{rrrrr}
\toprule
$N$ & $|S|$ & $|D|$ & prefixes & excluding prior job 24\\
\midrule
3&4&12&7&6\\
3&4&20&26&26\\
3&4&24&44&44\\
1&2&30&4&4\\
\midrule
&&&81&80\\
\bottomrule
\end{tabular}
\end{center}
The original job indices are
\begin{gather*}
24,26,28,29,31,33,34,36,38,39,40,41,42,43,44,45,46,\\
51,53,54,55,58,76,78,79,80,81,83,84,85,86,87,88,89,\\
91,92,93,95,96,98,100,101,102,104,106,108,109,111,112,\\
113,114,119,120,121,122,123,124,126,128,129,132,133,134,\\
135,136,137,138,139,141,143,144,145,148,149,150,159,\\
164,165,177,180,186.
\end{gather*}
Each index means the original rank-four job $[j,4]$; it is not a raw
involution index or a final-generator orbit position. The accompanying
application table pins the exact accepted source row, proof and result
for every listed prefix.

For job 24 the earlier proof already established $U=S\cong C_2^2$
and $|D|=12$, and applied the complement argument to exclude both
parabolics~\cite{old,oldcentralizer}. Its appearance in the table is
reuse, not another new restriction. The other 80 indices are distinct
applications relative to that explicitly retained job 24 record.
This is a statement about the present proof record, not a claim of
priority or absence from other work. The accompanying Root acceptance record adopts these conditional applications.

There are ten further records with $N=1$ and $|S|=2$, but their $D$
has order 66 or 110. They fail the required coprimality with 11 and
are not covered by this argument. The remaining 66 records fail
$N<5$. Ineligibility does not imply parabolic containment, and this
note supplies no converse or verdict for those 76 prefixes.

\section{Scope and attribution}

The use of orbit--stabilizer and Lagrange, coprime complement conjugacy,
the averaging proof and the finite split Levi structure are borrowed
standard mathematics. The earlier prefix note already assembled the
Levi-centre obstruction for job 24. The present assembly replaces the
unnecessary full equality $U=S$ by the proved inequality
$[U:S]\leq |\Inv(U)|$ and applies it to the exact admitted table.
No historical novelty of the abstract argument is asserted.

All model, membership, original frame, catalogue completeness and full
$S$ identifications are explicit inputs. If one of those imports is
not admitted for a proposed prefix, the corresponding application is
conditional until it is supplied. No imported premise is replaced by
an integer identity, a truncated scan or a mere known subgroup order.
For the present 81 records the adopted data supply the numerical
hypotheses; this note neither extends the full model proof nor fills
missing imports by assertion.

The output is a prefix restriction only. Excluding two maximal types
does not prove generation of a remaining end, the intersection property
or its failure for a quartet, a positive classification, a whole job or
a whole original manuscript. No final-generator position, paid union,
conjecture, Lean result or compression measurement is credited.
No Source, matrix, field, group, orbit, state, gzip or TeX execution
was performed while preparing this written application.

\begin{thebibliography}{9}
\bibitem{old}
Retained \emph{A fixed prefix excludes both parabolics}, corrected
publisher locators, SHA-256
\texttt{e9f31197a97fa2340a9ad5ca40c6d300f685d9dade0c35914486af40e6e7ce49}.
Its primary locator record has SHA-256 prefix \texttt{23d39eff23769888}.
\bibitem{Lin}
Z. Lin, \emph{The Lyndon--Demu\v{s}kin method and crystalline lifts of
$G_2$-valued Galois representations}, \emph{Algebra \& Number Theory}
\textbf{19} (2025), 415--456,
\href{https://doi.org/10.2140/ant.2025.19.415}{doi:10.2140/ant.2025.19.415},
Sections 7.1.4 and 7.1.6. Cited here through the retained primary
excerpts and finite-specialization record, not a fresh full-paper review.
\bibitem{facts}
Root-adopted complete involution centralizer facts for the current 157
prefixes, facts SHA-256
\texttt{2b7206d2dbcb171d5bf3df163e29cca71a531c731720f87ec9913954f28b3653};
acceptance prefix \texttt{a81814c2f2106318}. The recorded scan and exact
order certificates, including all original imports, are reused.
\bibitem{oldcentralizer}
Retained job 24 centralizer equality, SHA-256 prefix
\texttt{9d2f9c739a8a5f42}; its stronger equality is not a premise of
Corollary~\ref{q11}.
\end{thebibliography}
\end{document}
